\section{The local-gap engine}\label{sec:eng4}

\providecommand{\obs}{\operatorname{obs}}
\newcommand{\anout}{c_{\mathrm{out}}}
\newcommand{\anin}{c_{\mathrm{in}}}

The input of this section is an instance $\mathcal I$
(Definition~\ref{def:v4-instance}); the output is a rejection if $\mathcal I$
is infeasible, and otherwise a system of binary rows whose solutions,
restricted to $X$, are exactly $F(\mathcal I)$, with a witness for every
solution (Theorem~\ref{thm:eng4-main}). Witnesses, the contract, relocation
and the observers $\langle t\rangle$ are those of Section~\ref{sec:v4-lang}.
We use Lemma~\ref{lem:v4-gaps}(a),(b); the proof of (b) uses only
$\rho_{v'}=\rev\varphi(\rho_u)$ and that $\varphi$ maps the anchor pair at
$u$ onto that at $v'$, so we apply it also to such bijections that are not
pipes.

The engine runs the Synchronized Planarity algorithm of Bl\"asius, Fink and
Rutter~\cite{BFR21} (numbers refer to arXiv:2007.15362v2), with its graph
operations unchanged up to subdividing the edges at the leaves of the trees
inserted by PropagatePQ, and carries observations and rows along. Our pipe
condition $\rho_v=\rev\varphi(\rho_u)$ is that of~\cite{BFR21}. The
equivalence lemmas of~\cite{BFR21} concern feasibility of the bare instance,
whereas the contract preserves the feasible assignments of all existing
variables; so all statements about observations are proved here. We use
unchanged: the operations of~\cite[Sections~3.1.1--3.4]{BFR21}; splits and
joins along pipes~\cite[Section~2, Lemma~1]{BFR21}; the schedule in the
proof of~\cite[Lemma~7]{BFR21}; the potential of~\cite[Section~3.6]{BFR21} with~\cite[Lemma~9, Lemma~10(i), Lemma~11]{BFR21};
and that a wheel lies in one rigid skeleton~\cite[proof of Lemma~8]{BFR21}.
The relation to the rotation variables and triple equations
of~\cite[Chapter~8]{Pfr22} is discussed in the introduction.

\subsection{Wheels, embedding trees and Q-intersections}\label{sec:eng4-wheels}

The \emph{wheel graph} $G^W$ replaces every Q-vertex $q$ with $\deg q\ge3$
by a wheel~\cite[Section~2]{BFR21}: each edge $e$ at $q$ is subdivided by a
\emph{port} next to $q$ (twice if both ends are such Q-vertices), and the
ports are joined by a \emph{rim} in the order $\psi(q)$; $q$ is the
\emph{hub}. The neighbours of a hub are exactly its ports, pairwise distinct,
and its rim is a simple cycle. Darts of $G$ correspond to darts of $G^W$ at
the same vertex, and a marked cycle $C$ is \emph{lifted} to $C^W$ through the
spokes and ports of the hubs it passes, with the same anchors at P-vertices.

\begin{lemma}[wheels]\label{lem:eng4-wheel}
(a) Every planar $\rho$ of $G$ with $\rho_q\in\{\psi(q),\rev\psi(q)\}$ for
all hubs extends to a planar $\rho^W$ of $G^W$ with the same rotations at all
vertices of $G$. (b) Deleting the rims of a planar $\rho^W$ and suppressing
the ports gives a planar $\rho$ of $G$ with the same rotations at vertices of
$G$ and $\rho_q\in\{\psi(q),\rev\psi(q)\}$ at every hub. (c) A P-vertex is a
cutvertex of $G$ if and only if it is one of $G^W$. (d)
$\gap_\rho(C,v,d)=\gap_{\rho^W}(C^W,v,d)$ for every observable triple.
\end{lemma}

\begin{proof}
(a) Draw each rim on a small circle around its hub. (b) The wheel is $3$-connected, so the restriction of $\rho^W$ to it is one
of its two mirror embeddings, and the rotation at the hub is the rim order
$\psi(q)$ or its reverse. (c) Contracting the wheels of $G^W-w$ gives $G-w$. (d)
The rotation and the anchors at $v$ are the same.
\end{proof}

\paragraph{SPQR trees and embedding trees.} For a block $B$ of $G^W$ that is
neither an edge nor a cycle we use from Section~\ref{sec:v4-spqr}
\ref{it:v4-T1}--\ref{it:v4-T4}, reflection bits $r_\mu$, displayed cycles,
anchors and gaps $\gap_\mu$ of P skeletons (Lemma~\ref{lem:v4-display}),
embedding trees $\Theta_B(v)$ with frontiers, local reads and their complete
family (Lemmas~\ref{lem:v4-frontier}, \ref{lem:v4-trees}); and
Lemmas~\ref{lem:v4-median} and~\ref{lem:v4-assembly}. The \emph{embedding tree} $T_v$
of~\cite[Section~2]{BFR21} is $\Theta_B(v)$ with every P node of degree three
read as a Q node; it is \emph{trivial} if it is a single P node.

\begin{lemma}[Q-intersections]\label{lem:eng4-qint}
Let $q$ be a hub of $G^W$ and $B$ its block.
\begin{enumerate}[label=(\alph*),nosep]
\item $q$ is not a cutvertex of $G^W$.
\item $q$ is not a pole of a P node of $\mathcal T_B$.
\item Both poles of a P node with at least four skeleton edges are
  P-vertices of $G$.
\item Distinct lifted marked cycles meet only at hubs. Hence no cutvertex of
  $G^W$ lies on two lifted marked cycles, and at most one lifted marked cycle
  passes through both poles of a P node.
\end{enumerate}
\end{lemma}

\begin{proof}
(a) The rim connects all neighbours of $q$ in $G^W-q$. (b) Let $\mu$ have
poles $q,w$. The ports other than $w$ are inner to expansions, and the rim
minus $w$ is connected and avoids $q,w$; so they and their spokes lie in one
expansion, and at most the spoke $qw$ lies in another. So at most two
expansions contain darts at $q$, but $\mu$ has at least three skeleton
edges. (c) A pole $a$ has at least four darts, so it is not a Q-vertex of
degree at most two, nor a hub by (b). Let $a$ be a port of the hub $q$, with
rim neighbours $p',p''$; by (b) the other pole $b$ is not $q$. If
$b\notin\{p',p''\}$, the spoke and both rim edges at $a$ end in vertices
joined in $G^W-\{a,b\}$ through $q$; if $b=p'$, the spoke and $ap''$ end at
$q,p''$, joined by a spoke. Either way at most three expansions contain
darts at $a$, a contradiction. (d) Two marked cycles share only Q-vertices
and no edge; a Q-vertex of degree at most two or a port lies on at most one
lift. So shared vertices of lifts are hubs, neither cutvertices (a) nor
poles (b).
\end{proof}

\subsection{The operations}\label{sec:eng4-ops}

Each lemma below states that a step satisfies the contract and that its output is an instance, in
particular satisfies~\ref{it:v4-M1}--\ref{it:v4-M4}. Throughout,
$(u,v,\varphi)$ is the pipe acted on, $D=\deg u$, and observers are created
as $\langle t\rangle$. Old rows are kept, so it suffices to check planarity,
\ref{it:v4-W1}--\ref{it:v4-W3} and the new rows. Rotations change only at
the vertices acted on and, in SimplifyMatching, at the partner pole.



\begin{lemma}[ConvertSmall]\label{lem:eng4-convert}
Let $u$ be a P-vertex with $\deg u\le3$. Make $u$ a Q-vertex with any
reference $\psi(u)$ and a fresh flip $g=f(u)$. If $u$ is matched by
$(u,v,\varphi)$, delete the pipe and make $v$ a Q-vertex of the same Q-cell,
$f(v)=g$, with $\psi(v)=\rev\varphi(\psi(u))$. Relocate every observer of
$(C,z,d)$, $z\in\{u,v\}$, to $g\oplus\gap_{\psi(z)}(C,z,d)$. This step
satisfies the contract with $\rho'=\rho$ and yields an instance.
\end{lemma}

\begin{proof}
The invariant is clear. Observations at $z$ need $\deg z=3$. Forward: put
$\alpha'(g)=1$ if and only if $\deg u=3$ and $\rho_u=\rev\psi(u)$; then
$\rho_v=\rev\varphi(\rho_u)$ is $\psi(v)$ or $\rev\psi(v)$ accordingly,
and Lemma~\ref{lem:v4-gaps}(a) gives the new rows. Backward: for $\deg u=3$,
\ref{it:v4-W2} gives $(\rho_u,\rho_v)=(\psi(u),\rev\varphi\psi(u))$ or
$(\rev\psi(u),\varphi\psi(u))$, so the pipe holds (trivially if $\deg u\le2$),
and relocated observations hold by the same identity.
\end{proof}

\paragraph{Encapsulation.} Let $u$ be a P-vertex that is a cutvertex of its
component $W$, and $Y_1,\dots,Y_k$ the components of $W-u$.
As in~\cite[Section~3.2]{BFR21}, for $j=1,\dots,k$ split the component of
$u$ along the cut $(V(Y_j),\text{rest})$, creating the \emph{residual cap}
$u_j$ (contracted rest), the \emph{ray} $u'_j$ (contracted $Y_j$) and the cap
pipe $(u_j,u'_j,\chi_j)$. An edge $e=uy$, $y\in Y_j$, becomes a \emph{star
edge} $uu'_j$, which keeps the dart of $e$ at $u$, and a \emph{residual edge}
$u_jy$, which keeps the dart at $y$; so old vertices keep their dart sets and
rotations. The component of $u$ becomes the star $S_u$. If $u\in
C\in\mathcal M$, the anchors of $u$ end in one $Y_{j_0}$; with $\hat
u=u'_{j_0}$, $C$ is replaced by $C^{\mathrm{res}}$ through $u_{j_0}$ and
the residual edges, and the \emph{star digon} $C^\star_u:u\to\hat u\to u$
(out along the star edge of $\anout(C,u)$, back along that of $\anin(C,u)$)
is added. Observations at $u$ now refer to $C^\star_u$, with the same dart;
those at other vertices of $C$ to $C^{\mathrm{res}}$.

\begin{lemma}[encapsulation]\label{lem:eng4-encap}
The encapsulation satisfies the contract with rotations unchanged at all old
vertices, and its result is an instance. A vertex $y\ne u$ of $W$ is a
cutvertex of $W$ if and only if it is one of the piece $Y_j+u_j$
containing it, and every residual cap is a block-vertex of its piece.
\end{lemma}

\begin{proof}
Both shores of each cut are connected, so by~\cite[Lemma~1, Section~2]{BFR21}
every planar rotation system splits into a planar one with the same old
rotations satisfying the cap pipes, and conversely the joins along the cap
pipes are planar and keep old rotations. Observations keep their values,
as rotations and anchors are unchanged.
\emph{Invariant.} $C^{\mathrm{res}}$ and $C^\star_u$ are simple, the new
vertices on them lie on one marked cycle each, $u$ lies only on
$C^\star_u$, and other marked cycles are unchanged. The cap pipe at $j_0$
maps anchors to anchors, hence is aligned; a marked cycle through $u_j$ or
$u'_j$, $j\ne j_0$, would cross the $j$-th cut, so the other cap pipes are
off. \emph{Cutvertices.} The components of $W-y$ avoiding $u$ lie in $Y_j$
and survive in $(Y_j+u_j)-y$, and $u_j$ lies in another one. Conversely, if
$W-y$ is connected, every component of $Y_j-y$ is adjacent to $u$, hence to
$u_j$. Finally $(Y_j+u_j)-u_j=Y_j$ is connected.
\end{proof}

\paragraph{Joining two stars.}
Let $S_u,S_v$ be disjoint stars (parallel edges allowed) with centres $u,v$
and ray sets $R_u,R_v$, and let $\varphi:A_u\to A_v$ be a bijection. For
$d\in A_u$ let $r(d)\in R_u$ be the ray at the other end of $d$, and
$r'(d')\in R_v$ likewise for $d'\in A_v$. The
\emph{join}~\cite[Section~2]{BFR21} deletes $u$ and $v$ and adds, for every
$d\in A_u$, an edge $\ell_d=r(d)\,r'(\varphi d)$ that takes over the ray
darts of $d$ and $\varphi(d)$. The result is the bipartite \emph{ray graph}
$K$; a rotation of a ray is read in $S_u\cup S_v$ and in $K$ alike, darts
renamed.

\begin{lemma}[stars]\label{lem:eng4-stars}
\begin{enumerate}[label=(\alph*),nosep]
\item A rotation system of two vertices joined by $k\ge1$ parallel edges is
  planar if and only if the rotation of one vertex is the reverse of that of
  the other, edges identified.
\item Let $\sigma$ be a cyclic order of $A_u$ and $\rho'$ rotations of the
  rays. If the rotation system of $S_u\cup S_v$ given by $\sigma$ at $u$,
  $\rev\varphi(\sigma)$ at $v$ and $\rho'$ at the rays is planar, then so
  is $\rho'$ on $K$.
\item Conversely, if $\rho'$ is a planar rotation system of $K$ and $K_j$ is
  a component of $K$, there is a cyclic order $\sigma_j$ of
  $\{d:\ell_d\in K_j\}$ such that the stars $S^j_u,S^j_v$ spanned by these
  darts and by their images, with centre orders $\sigma_j$,
  $\rev\varphi(\sigma_j)$ and $\rho'$ at the rays, are planar; it is
  computable in $O(|E(K_j)|)$ time.
\end{enumerate}
\end{lemma}

\begin{proof}
(a) Read the rotations $\sigma,\pi$ as cyclic permutations of the $k$
edges. A facial walk leaving the first vertex along $e$ leaves it next along
$\sigma\pi(e)$, so facial walks correspond to cycles of $\sigma\pi$. By
Euler's formula the system is planar if and only if there are $k$ of them,
that is, $\sigma\pi=\mathrm{id}$, or $\pi=\rev\sigma$.

(b), (c) The blocks of a star are the bonds between its centre and its rays.
By (a) and Lemma~\ref{lem:v4-assembly}(a),(c), the rotation system of
$S_u\cup S_v$ described in (b) is planar if and only if: every ray $x\in
R_u$ has the rotation $\rev(\sigma|_x)$, where $\sigma|_x$ is the
restriction of $\sigma$ to the darts ending at $x$; every ray $y\in R_v$ has
the rotation $\rev\bigl(\rev\varphi(\sigma)|_y\bigr)$, which is $\sigma|_y$
after identifying the dart of $\ell_d$ at $y$ with $d$; and the partitions
of $A_u$ by $r$ and by $r'\circ\varphi$ are noncrossing in $\sigma$. Put
$\omega=\rev\sigma$. The condition says that $\omega$ restricts to
$\rho'_x$ at every $x\in R_u$ and to $\rev\rho'_y$ at every $y\in R_v$, and
makes both partitions noncrossing. The multigraph of
Lemma~\ref{lem:v4-partitions} for these partitions of $A_u$, with
$\sigma_X=\rho'_x$ at the red blocks $x\in R_u$ and $\sigma_X=\rev\rho'_y$
at the blue blocks $y\in R_v$, is $K$ with the rotation system $\rho'$. So
Lemma~\ref{lem:v4-partitions} gives (b), and, applied to the darts of one
component, (c).
\end{proof}

\paragraph{Gap labels.}
Let $\rho$ be a planar rotation system of $S_u\cup S_v$ with
$\rho_v=\rev\varphi(\rho_u)$, and $\rho'$ its rotations at the rays, which
are planar on $K$ by Lemma~\ref{lem:eng4-stars}(b). Suppose that $a\ne b$
are darts at $u$ with $r(a)=r(b)=:p$ and $r'(\varphi a)=r'(\varphi b)=:q$.
Let $C_u:u\to p\to u$ be the digon that leaves $u$ along $a$ and returns
along $b$, and let
\[
C^\ast:\ p\xrightarrow{\ \ell_b\ }q\xrightarrow{\ \ell_a\ }p .
\]
The \emph{label} of an edge $\ell_d$, $d\ne a,b$, is
$h_d=\gap_\rho(C_u,u,d)$. A ray other than $p$ and $q$ is \emph{ordinary};
all edges at an ordinary ray are labelled. Let $K_0$ be the component of $K$
containing $C^\ast$. For a dart $d\ne a,b$ at $u$ with $\ell_d\in K_0$ define
the \emph{target} $t(d)$: it is $(C^\ast,p,\ell_d)$ if $r(d)=p$; otherwise
$(C^\ast,q,\ell_d)$ if $r'(\varphi d)=q$; otherwise fix a path of $K_0$ from
$r(d)$ that meets $\{p,q\}$ only in its last vertex $z$, let $g$ be its last
edge, and put $t(d)=(C^\ast,z,g)$.

\begin{lemma}[gap labels]\label{lem:eng4-labels}
In this situation:
\begin{enumerate}[label=(\alph*),nosep]
\item the edges at an ordinary ray have one label, as have the edges of a
  component $T\ne K_0$;
\item $\gap_{\rho'}(C^\ast,z,\ell_d)=h_d$ for $z\in\{p,q\}$ and every edge
  $\ell_d\ne\ell_a,\ell_b$ at $z$;
\item every target is an observable triple of $C^\ast$, and
  $\gap_{\rho'}(t(d))=h_d$ for every $d\ne a,b$ with $\ell_d\in K_0$.
\end{enumerate}
\end{lemma}

\begin{proof}
(a) Let $x\in R_u\setminus\{p\}$. The blocks of $S_u$ are the bonds between
$u$ and its rays, so by Lemma~\ref{lem:v4-assembly}(a) the darts at $u$ of
$x$ do not interleave with $a,b$, which belong to $p$. Hence they all lie in
$(a,b)_{\rho_u}$ or all in $(b,a)_{\rho_u}$, and their labels agree. For
$x\in R_v\setminus\{q\}$, likewise the darts at $v$ of $x$ lie all in
$(\varphi a,\varphi b)_{\rho_v}=\varphi((b,a)_{\rho_u})$ or all in
$(\varphi b,\varphi a)_{\rho_v}=\varphi((a,b)_{\rho_u})$, as
$\rho_v=\rev\varphi(\rho_u)$, and the labels agree. A component $T\ne K_0$ avoids
$p$ and $q$, so all its vertices are ordinary, and by connectivity all its
edges have one label.

(b) By Lemma~\ref{lem:eng4-stars}(a), applied to the bond between $u$ and
$p$, the rotation of $p$ is the reverse of the restriction $\sigma$ of
$\rho_u$ to the darts ending at $p$. Now $C^\ast$ leaves $p$ along $\ell_b$
and enters it along $\ell_a$, so for $\ell_d$ at $p$, $d\ne a,b$:
$\ell_d\in(\ell_b,\ell_a)_{\rho'_p}$ if and only if
$d\in(b,a)_{\rev\sigma}=(a,b)_\sigma$, if and only if $d\in(a,b)_{\rho_u}$.
Thus $\gap_{\rho'}(C^\ast,p,\ell_d)=h_d$. At $q$, where $C^\ast$ enters along
$\ell_b$ and leaves along $\ell_a$, symmetrically
$\ell_d\in(\ell_a,\ell_b)_{\rho'_q}$ if and only if $\varphi
d\in(\varphi b,\varphi a)_{\rho_v}=\varphi((a,b)_{\rho_u})$.

(c) If $r(d)=p$ or $r'(\varphi d)=q$, the claim is (b) (and $\ell_d\ne
\ell_a,\ell_b$). Otherwise $r(d)$ is ordinary, the path defining $t(d)$
exists as $K_0$ is connected, and all its vertices except $z$ are ordinary.
Its last edge $g$ has an ordinary end, so $g\ne\ell_a,\ell_b$ and $t(d)$ is
observable. By (a), $\ell_d$ and the first edge of the path have the same
label, and so have any two consecutive edges of the path; so $g$ has label
$h_d$, and (b) gives $\gap_{\rho'}(C^\ast,z,g)=h_d$.
\end{proof}

\begin{figure}[t]
\centering
\begin{tikzpicture}[
  every node/.style={font=\small},
  uv/.style={circle,fill=black,inner sep=1.4pt},
  vv/.style={circle,draw=black,fill=white,inner sep=1.4pt},
  ctr/.style={circle,draw=black,fill=black!15,inner sep=1.8pt},
  >={Stealth[length=2.2mm]}]
\begin{scope}
\node[ctr,label=below:$u$] (u) at (0,0) {};
\node[uv,label=above:$p$] (p) at (-0.6,2.2) {};
\node[uv,label={[label distance=-1pt]left:$t$}] (t) at (-0.25,1.0) {};
\node[uv,label=left:$x$] (x) at (-1.6,-1.1) {};
\draw[->,thick] (u) to[out=60,in=-40] node[pos=0.45,right]{$a$} (p);
\draw[->,thick] (p) to[out=250,in=135] node[pos=0.3,left=-1pt]{$b$} (u);
\draw (u) to[out=170,in=200] node[pos=0.5,left]{\scriptsize$d_1$} (p);
\draw (u) -- node[pos=0.45,right=-1pt]{\scriptsize$d_3$} (t);
\draw (u) to[bend right=15] node[pos=0.5,above left=-2pt]{\scriptsize$d_2$} (x);
\draw (u) to[bend left=15] node[pos=0.5,below right=-2pt]{\scriptsize$d_4$} (x);
\node[ctr,label=below:$v$] (v) at (3.4,0) {};
\node[vv,label=above:$q$] (q) at (4.2,2.2) {};
\node[vv,label={[label distance=-1pt]right:$y$}] (y) at (3.55,1.1) {};
\node[vv,label=right:$t'$] (tt) at (2.9,-1.3) {};
\draw[thick] (v) to[out=15,in=-10] (q);
\draw (v) to[out=50,in=-80] (q);
\draw[thick] (v) to[out=125,in=170] (q);
\draw (v) to[bend right=22] (y);
\draw (v) to[bend left=22] (y);
\draw (v) -- (tt);
\draw[dashed] (u) -- node[below]{$\varphi$} (v);
\node at (1.7,-1.9) {(a) stars $S_u$, $S_v$ and the pipe};
\end{scope}
\begin{scope}[xshift=7.6cm]
\node[uv,label=left:$p$] (P) at (0,0.2) {};
\node[vv,label=right:$q$] (Q) at (3,0.2) {};
\node[vv,label=above:$y$] (Y) at (0.8,1.7) {};
\node[uv,label=above:$x$] (X) at (2.2,1.7) {};
\draw[->,thick] (P) to[bend right=25] node[pos=0.5,below,inner sep=1pt]{$\ell_b$} (Q);
\draw[->,thick] (Q) to[bend right=25] node[pos=0.5,above,inner sep=1pt]{$\ell_a$} (P);
\draw (P) -- node[left,inner sep=1pt]{$\ell_{d_1}$} (Y);
\draw (Y) -- node[above,inner sep=1pt]{$\ell_{d_2}$} (X);
\draw (X) -- node[right,inner sep=1pt]{$\ell_{d_4}$} (Q);
\node at (1.5,0.2) {$C^\ast$};
\node[uv,label=left:$t$] (T1) at (1.0,-1.1) {};
\node[vv,label=right:$t'$] (T2) at (2.0,-1.1) {};
\draw (T1) -- node[below,inner sep=1pt]{$\ell_{d_3}$} (T2);
\node at (1.5,-1.9) {(b) the ray graph $K$};
\end{scope}
\end{tikzpicture}
\caption{The join of two stars; filled rays belong to $u$, hollow rays to
$v$, and the drawings realize the rotations. Counter-clockwise around $u$
the darts are $a,d_3,b,d_1,d_2,d_4$. The anchor edges $\ell_a,\ell_b$ form
the digon $C^\ast$. The component $K_0$ contains the bridge of $C^\ast$
through $y$ and $x$, whose edges all have label $1$; this is the local gap
of $\ell_{d_1}$ at $p$ and of $\ell_{d_4}$ at $q$. The detached component
$T=\{tt'\}$ receives the latent bit $h_T=h_{d_3}=0$.}
\label{fig:eng4-join}
\end{figure}

\begin{definition}[EncapsulateAndJoin]\label{def:eng4-ej}
Let $(u,v,\varphi)$ be a pipe whose ends are cutvertices of $G$ (possibly in
one component, possibly adjacent). Encapsulate $u$, then $v$; call the
result $\mathcal I^\ast$, with stars $S_u,S_v$. \emph{Join}
$S_u$ and $S_v$ along the pipe~\cite[Section~3.2]{BFR21}: delete $u$, $v$,
their edges and the pipe, and add the ray graph $K$; the rays keep their cap
pipes. If $K$ is not planar, reject. If the pipe is aligned ($u\in C_u$,
$v\in C_v$), let $a=\anout(C_u,u)$, $b=\anin(C_u,u)$, $p=\hat u$, $q=\hat v$;
replace $C^\star_u$ and $C^\star_v$ by the marked digon $C^\ast:
p\xrightarrow{\ell_b}q\xrightarrow{\ell_a}p$, add a fresh \emph{latent bit}
$h_T$ for every component $T\ne K_0$ of $K$, and relocate every observer of
$(C^\star_u,u,d)$ to $\langle t(d)\rangle\oplus0$ if $\ell_d\in K_0$ and to
$h_T\oplus0$ if $\ell_d\in T$, and every observer of
$(C^\star_v,v,\varphi(d))$ in the same way with offset
$\varepsilon(\varphi)$ in place of $0$.
\end{definition}

\begin{lemma}[EncapsulateAndJoin]\label{lem:eng4-join}
EncapsulateAndJoin satisfies the contract, yields an
instance, and rejects only if $K$ is not planar. It adds $O(D)$ vertices,
edges, pipe incidences, variables and rows.
\end{lemma}

\begin{proof}
By Lemma~\ref{lem:eng4-encap} it suffices to treat $\mathcal
I^\ast\to\mathcal I'$; before the second encapsulation $v$ is still a
cutvertex (of its piece), in the aligned case on the image of $C_v$ with
the same anchors. In
$\mathcal I^\ast$ every observation in the stars is at $u$ or $v$, and in the
aligned case $r(a)=r(b)=p$ and $r'(\varphi a)=r'(\varphi b)=q$, so
Lemma~\ref{lem:eng4-labels} applies with $C_u=C^\star_u$.
\emph{Off pipe.} No observation lies at $u$ or $v$, no marked cycle meets
the stars, and the step is the join of~\cite{BFR21}. Its inverse is the
backward construction below with any component as $K_0$ and any dart $a$ of
it, all other components inserted after $a$; both keep rotations and
anchors at observed vertices. From now on the pipe is aligned.
\emph{Invariant.} $K$ is bipartite and loopless, rays keep their degrees,
$C^\ast$ is a digon on new edges, no other marked cycle meets a ray, the cap
pipes at $p,q$ stay aligned and the others off; targets are observable
(Lemma~\ref{lem:eng4-labels}(c)).
\emph{Forward.} Given a witness $(\rho,\alpha)$ of $\mathcal I^\ast$, let
$\rho'$ consist of its rotations at the rays and $\rho$ elsewhere; it is
planar by Lemma~\ref{lem:eng4-stars}(b), and the cap pipes hold. Give new
observers their gaps and $h_T$ the common label of $T$
(Lemma~\ref{lem:eng4-labels}(a)). The observation of $(C^\star_u,u,d)$ has
value $h_d$, which is $\alpha'(h_T)$ if $\ell_d\in T$ and
$\gap_{\rho'}(t(d))$ if $\ell_d\in K_0$ (Lemma~\ref{lem:eng4-labels}(c)); at
$v$, Lemma~\ref{lem:v4-gaps}(b) gives
$\gap_\rho(C^\star_v,v,\varphi d)=h_d\oplus\varepsilon(\varphi)$.

\emph{Backward.} Given a witness $(\rho',\alpha')$ of $\mathcal I'$,
Lemma~\ref{lem:eng4-stars}(c) gives for every component $K_j$ planar stars
$S^j_u,S^j_v$ with the ray rotations of $\rho'$ and centre orders $\omega_j$,
$\rev\varphi(\omega_j)$. Start with $\rho_u=\omega_0$,
$\rho_v=\rev\varphi(\omega_0)$. For every $T=K_j$, $j\ge1$, insert a linear
order $L_T$ representing $\omega_j$ into $\rho_u$ directly after $a$ if
$\alpha'(h_T)=0$ and directly after $b$ if $\alpha'(h_T)=1$, and
$\rev\varphi(L_T)$ into $\rho_v$ directly before $\varphi(a)$, respectively
$\varphi(b)$; this preserves $\rho_v=\rev\varphi(\rho_u)$. The final
$\rho_u$ restricts to every $\omega_j$, and dart sets of distinct $K_j$ do
not interleave, as each insertion is contiguous. So the bonds of $S_u$ are
planar and pairwise noninterleaving, and $S_u$ is planar by
Lemma~\ref{lem:v4-assembly}(c); likewise $S_v$; the pipes hold. Later
insertions after $a$ or $b$ keep earlier blocks in their interval, so
$\gap_\rho(C^\star_u,u,d)=\alpha'(h_T)$ for $\ell_d\in T$; and they keep
the cyclic order of $\omega_0$, so for $\ell_d\in K_0$,
$\gap_\rho(C^\star_u,u,d)=\gap_{\omega_0}(C^\star_u,u,d)=\gap_{\rho'}(t(d))$
by Lemma~\ref{lem:eng4-labels}(c) for $S^0_u,S^0_v$, whose join is $K_0$.
Observations at $v$ follow by Lemma~\ref{lem:v4-gaps}(b).
\emph{Rejection and size.} The forward construction does not use rows, so a
witness forces $K$ to be planar. The step adds $O(D)$ rays, cap pipes, edges
and targets, at most $D-2$ latent bits and at most $2(D-2)$ relocations.
\end{proof}

\paragraph{PropagatePQ.} Let $u$ be a matched block-vertex whose embedding
tree $T_u$, computed in the block of $G^W$ containing $u$, is not trivial. As
in~\cite[Section~3.3]{BFR21}, replace $u$ by $T_u$: its leaves become
vertices $\ell_d$ adjacent to their inner node and to the other end of $d$
(whose dart there keeps its identity); replace $v$ by a copy $T'_u$ under an
isomorphism $\kappa$, the leaf $\kappa(\ell_d)$ receiving the edge of
$\varphi(d)$ (an edge $uv$ thus joins two leaves); for each inner node
$\alpha$ let $\kappa_\alpha$ map the darts at $\alpha$ to those at
$\kappa\alpha$ (tree darts to tree darts). Add the pipe
$(\alpha,\kappa\alpha,\kappa_\alpha)$ for every P node $\alpha$. For every Q
node $\alpha$ make $\alpha$ and $\kappa\alpha$ Q-vertices of one fresh
Q-cell with a fresh flip $g_\alpha=f(\alpha)=f(\kappa\alpha)$,
$\psi(\alpha)=\theta_\alpha$, the reference order of the
node (either order for a P node of degree three read as a Q node), and
$\psi(\kappa\alpha)=\rev\kappa_\alpha(\theta_\alpha)$. Delete the pipe
$(u,v,\varphi)$. If $u\in C$, replace the visit of $C$ at $u$ by the path of
$T_u$ from $\ell_{\anin(C,u)}$ to $\ell_{\anout(C,u)}$ (the new cycle $C^T$;
observations at other vertices of $C$ now refer to $C^T$, with unchanged
anchors); treat $v$ and $C(v)$ likewise in $T'_u$. \emph{Median transport:}
for an observer $x$ of $(C,u,d)$ let $m$ be the median of
$\ell_{\anout(C,u)},\ell_{\anin(C,u)},\ell_d$ in $T_u$ and $\delta$ the dart
of $m$ towards $\ell_d$; relocate $x$ to $\langle C^T,m,\delta\rangle\oplus0$
if $m$ is a P node and to $f(m)\oplus\gap_{\theta_m}(C^T,m,\delta)$ if $m$ is
a Q node. Observers at $v$ are treated in the same way in $T'_u$.

\begin{lemma}[PropagatePQ]\label{lem:eng4-ppq}
This step satisfies the contract and yields an instance. It adds $O(D)$
vertices, edges, pipe incidences, variables and rows.
\end{lemma}

\begin{proof}
\emph{Invariant.} $C^T$ replaces $u$ by a path of new vertices, which lie on
no other marked cycle. As $\varphi$ maps the anchors of $u$ onto those of
$v$, the path of $C(v)^T$ in $T'_u$ is $\kappa$ of the path $P$ of $C^T$, so
the new pipes are aligned on $P$ and off elsewhere. The median $m$ lies on
$P$, $\delta$ does not, and $C^T$ enters $m$ towards $\ell_{\anin(C,u)}$ and
leaves towards $\ell_{\anout(C,u)}$, so targets are observable.
\emph{Median identity.} If $\rho'$ is planar on $G'$, contracting $T_u$ and
$T'_u$ gives a planar $\rho$ with the frontiers at $u,v$
(Lemma~\ref{lem:v4-median}(d)), and by Lemma~\ref{lem:v4-median}(e)
$\gap_\rho(C,u,d)=\gap_{\rho'}(C^T,m,\delta)$, which is
$f(m)\oplus\gap_{\theta_m}(C^T,m,\delta)$ at a Q node by~\ref{it:v4-W2} and
Lemma~\ref{lem:v4-gaps}(a); likewise at $v$.
\emph{Forward.} By Lemmas~\ref{lem:eng4-wheel}(a)
and~\ref{lem:v4-frontier}(b), $\rho_u$ is the frontier of $T_u$ with the
rotations induced by an embedding of the block. Give $T'_u$ the mirrored
rotations $\rev\kappa_\alpha(\rho'_\alpha)$; its frontier is
$\rev\varphi(\rho_u)=\rho_v$~\cite[proof of Lemma~5]{BFR21}, so by
Lemma~\ref{lem:v4-median}(d) the expansion $\rho'$ is planar; the new pipes
hold, and $\alpha'(g_\alpha)=0$ if and only if $\rho'_\alpha=\theta_\alpha$
satisfies~\ref{it:v4-W2} at $\alpha$ and $\kappa\alpha$. The median
identity gives the rows. \emph{Backward.} Contract the trees; pipes and
Q-cells make them mirror images, so $\rho_v=\rev\varphi(\rho_u)$~\cite[proof
of Lemma~5]{BFR21}, and the median identity gives the observations at $u,v$.
\emph{Size.} $T_u$ has $\le D-2$ inner nodes and $D$ leaves.
\end{proof}

\paragraph{SimplifyMatching: the trivial pole.} In the rest of this
subsection $G^W$ is planar, every P-vertex has degree at least four, $u$ is a
matched block-vertex of $G$ with trivial embedding tree, $D=\deg u$, and $B$
is the block of $G^W$ containing $u$. As $T_u$ is trivial, $\Theta_B(u)$
consists of one P node and $D$ leaves: exactly one P or R node $\mu$ of
$\mathcal T_B$ allocates $u$, it is a P node with $D$ skeleton edges
$\varepsilon_1,\dots,\varepsilon_D$, and the expansion $B_i$ of
$\varepsilon_i$ contains exactly one edge at $u$; let $e_i$ be its dart at
$u$. The other pole $w$ of $\mu$ is the \emph{partner} of $u$. Let $I_i$ be
the set of darts at $w$ of the edges of $B_i$.

\begin{lemma}[trivial pole]\label{lem:eng4-pole}
\begin{enumerate}[label=(\alph*),nosep]
\item $w$ is a P-vertex of $G$; the sets $I_i$ are nonempty and partition
  the darts of $B$ at $w$; so $\deg w\ge D$.
\item Every directed simple cycle of $G^W$ through $w$ but not through $u$
  has its two darts at $w$ in one $I_i$ or in one block other than $B$.
\item If $u\in C\in\mathcal M$, then $w\in C$; the anchors
  $e_a=\anout(C,u)$, $e_b=\anin(C,u)$ satisfy $a\ne b$; $C^W$ is a path of
  $B_a$ from $u$ to $w$ followed by a path of $B_b$ from $w$ to $u$; and
  $\anin(C,w)\in I_a$, $\anout(C,w)\in I_b$.
\item If $\deg w=D$, or if $w$ is a block-vertex of $G$ with trivial
  embedding tree, then $\deg w=D$, every $I_i$ has exactly one element, and
  all darts at $w$ lie in $B$.
\end{enumerate}
\end{lemma}

\begin{proof}
(a) Lemma~\ref{lem:eng4-qint}(c), as $D\ge4$; each expansion is connected
and contains $w$. (b) A cycle lies in one block. If $Z\subseteq B-u$ met
two expansions, a maximal subpath of $Z$ in one of them would have both ends
at $w$, so it would be $Z$. (c) $C^W$ lies in $B$, and its darts at $u$ lie
in distinct expansions. Leaving $u$ along $e_a$ it stays in $B_a$ until it
reaches $w$, and symmetrically it returns through $B_b$; it passes $w$ once.
(d) If $\deg w=D$, the $D$ nonempty disjoint sets $I_i$ exhaust $A_w$. If
$w$ is a block-vertex with trivial embedding tree, $\Theta_B(w)$ is one P
node, which must be $\mu$, with one leaf per direction.
\end{proof}

Let $\rho$ be a planar rotation system of $G$ with
$\rho_q\in\{\psi(q),\rev\psi(q)\}$ for every Q-vertex of degree at least
three, $\rho^W$ its extension (Lemma~\ref{lem:eng4-wheel}(a)), and
$\mathcal E$ the embedding of $B$ induced by $\rho^W$. In $\mathcal E(w)$ the
sets $I_i$ form intervals appearing in the order $\mathcal E_\mu(w)$, and in
$\mathcal E(u)=\rho_u$ the darts $e_i$ appear in the order $\mathcal
E_\mu(u)=\rev\mathcal E_\mu(w)$ by \ref{it:v4-T3}, \ref{it:v4-T4}; so $I_k$
directly follows $I_i$ in $\mathcal E(w)$ if and only if $e_i$ directly
follows $e_k$ in $\rho_u$. Each $I_i$ is an interval of
$\mathcal E(w)=\rho_w|_{B}$ and so has a first dart. The \emph{superblock}
$S_i$ is the interval of $\rho_w$ from the first dart of $I_i$ up to,
excluding, the first dart of the next interval $I_k$, with its linear order.
The superblocks partition $A_w$, $I_i\subseteq S_i$, and the other darts of
$S_i$ belong to blocks other than $B$ (Figure~\ref{fig:eng4-superblock}).

\begin{figure}[t]
\centering
\def\sbToW#1#2{\draw (#1) to[out=0,in=#2] ($(w)+(#2:1.3)$) -- (w);}
\begin{tikzpicture}[scale=0.75,
  every node/.style={font=\small},
  pole/.style={circle,draw=black,fill=black!15,inner sep=1.8pt},
  bnd/.style={rectangle,rounded corners=2.2mm,draw=black,fill=black!8,
    inner sep=0pt,minimum width=8mm,minimum height=4.4mm,font=\scriptsize},
  oth/.style={draw=black,fill=black!40,line join=round},
  markcyc/.style={draw=black!18,line width=6pt,line cap=round,line join=round},
  >={Stealth[length=2.2mm]}]
\foreach \x/\l/\m/\pa/\pb/\pc/\ps/\pe/\cp/\cw in
  {0/2/3/124/140/156/116/170/{(a) $\rho_u=(e_5,e_4,e_3,e_2,e_1)$}/{$S_1S_2S_3S_4S_5$},
   7.6/3/2/146/162/178/138/192/{(b) $\tau=(e_5,e_4,e_2,e_3,e_1)$}/{$S_1S_3S_2S_4S_5$}} {
\begin{scope}[xshift=\x cm]
\coordinate (u) at (0,0); \coordinate (w) at (4.9,0);
\draw[markcyc] (u) to[out=55,in=180] (1.5,1.5) -- (2.1,1.5) to[out=0,in=100]
  ($(w)+(100:1.3)$) -- (w) -- ($(w)+(208:1.3)$) to[out=208,in=0] (2.1,-0.75)
  -- (1.5,-0.75) to[out=180,in=-27] (u);
\path (1.8,1.5) node[bnd] (b1) {$B_1$} (1.8,0.75) node[bnd] (b2) {$B_\l$}
  (1.8,0) node[bnd] (b3) {$B_\m$} (1.8,-0.75) node[bnd] (b4) {$B_4$}
  (1.8,-1.5) node[bnd] (b5) {$B_5$};
\foreach \b/\a in {b1/55,b2/27,b3/0,b4/-27,b5/-55}
  \draw (u) to[out=\a,in=180] (\b.west);
\sbToW{b1.east}{100} \sbToW{b4.east}{208} \sbToW{b5.east}{240}
\ifnum\l=2 \sbToW{$(b2.east)+(0,0.07)$}{\pa} \sbToW{$(b2.east)+(0,-0.07)$}{\pb}
  \sbToW{b3.east}{180}
\else \sbToW{b2.east}{124}
  \sbToW{$(b3.east)+(0,0.07)$}{\pa} \sbToW{$(b3.east)+(0,-0.07)$}{\pb}
\fi
\filldraw[oth] (w) -- ($(w)+(\pc:1.15)$) -- ($(w)+(\pc+11:1.15)$) -- cycle;
\node[font=\scriptsize,inner sep=1pt,anchor=east] at ($(w)+(\pc+5:1.18)$) {$\beta$};
\filldraw[oth] (w) -- ++(18:0.75) -- ++(270:0.46) -- cycle;
\node[font=\scriptsize,anchor=west] at ($(w)+(0:0.72)$) {$\beta'$};
\draw[densely dotted,thick] ($(w)+(\ps:0.62)$) arc[start angle=\ps,end angle=\pe,radius=0.62];
\draw[->,black!55] (0.55,1.25) -- (0.95,1.5);
\node at (0.45,1.6) {$C$};
\node[pole,label=left:$u$] at (u) {};
\node[pole,label={[label distance=1pt]below right:$w$}] at (w) {};
\node at (2.5,-2.25) {\cp};
\node at (2.5,-2.75) {at $w$: \cw};
\end{scope}}
\end{tikzpicture}
\caption{A trivial pole $u$, its partner $w$ and the five skeleton edges of
$\mu$, drawn as their expansions $B_1,\dots,B_5$; $e_i$ is the dart of
$B_i$ at $u$, the marked cycle $C$ (shaded) runs through $B_1$ and $B_4$,
$\beta,\beta'$ are blocks at $w$ other than $B$, and the drawings realize
the rotations. The superblock $S_2$ (dotted) consists of the darts of $B_2$
and $\beta$ at $w$, and $S_5$ of those of $B_5$ and $\beta'$. Exchanging
$e_2,e_3$ changes only the skeleton embedding of $\mu$ and moves $S_2,S_3$
as wholes; as $e_2,e_3$ stay in $(e_4,e_1)$ at $u$, all gaps at $u$ and $w$
are kept (Lemma~\ref{lem:eng4-superblock}(c)).}
\label{fig:eng4-superblock}
\end{figure}

\begin{lemma}[superblock permutation]\label{lem:eng4-superblock}\label{cor:eng4-fgb}
Let $\rho$ be as above and $\tau$ any cyclic order of $A_u$. Define
$\rho^\tau$ by $\rho^\tau_u=\tau$, $\rho^\tau_z=\rho_z$ for $z\ne u,w$, and
$\rho^\tau_w$ the cyclic concatenation of the superblocks, each with its
linear order, in which $S_k$ directly follows $S_i$ if and only if $e_i$
directly follows $e_k$ in $\tau$.
\begin{enumerate}[label=(\alph*),nosep]
\item $\rho^\tau$ is planar, and $\rho^{\rho_u}=\rho$.
\item If $u$ lies on no marked cycle, every observation has the same value
  under $\rho$ and $\rho^\tau$.
\item If $u\in C\in\mathcal M$ and $\gap_\tau(C,u,e)=\gap_\rho(C,u,e)$ for
  every nonanchor $e$ at $u$, every observation has the same value under
  $\rho$ and $\rho^\tau$.
\end{enumerate}
\end{lemma}

\begin{proof}
(a) $\rho^{\rho_u}=\rho$ was shown above. Let $\mathcal E^\tau$ be the
embedding of $B$ given by \ref{it:v4-T3} with the skeleton embeddings of $\mathcal
E$ except $\mathcal E^\tau_\mu(u)=\tau$. Its rotations are the frontiers of
the embedding trees with the induced node rotations. For $x\in
V(B)\setminus\{u,w\}$, $\mu$ does not allocate $x$, so $\mathcal
E^\tau(x)=\mathcal E(x)$; and $\mathcal E^\tau(u)=\tau$. In $\Theta_B(w)$
only the rotation of $\mu$ changes, to $\rev\tau$; by
Lemma~\ref{lem:v4-median}(b) the contour concatenates the excursions at
$\mu$ in this order, each visiting its leaves in an unchanged order. So
$\mathcal E^\tau(w)$ concatenates the same linear sequences $I_i$, with $I_k$
directly after $I_i$ if and only if $e_i$ directly follows $e_k$ in $\tau$;
this is the restriction of $\rho^\tau_w$ to $B$.

Now apply block assembly (Lemma~\ref{lem:v4-assembly}(c)) to the component
of $u$ in $G^W$, with $\mathcal E^\tau$ on $B$, the restrictions of $\rho^W$
on the other blocks, $\rho^W_x$ at cutvertices $x\ne w$ (where nothing
changed; $u$ is no cutvertex) and $\rho^\tau_w$ at $w$. Every other block
$B'$ at $w$ does not interleave with $B$ in $\rho_w$, so its darts lie in
one \emph{slot} between consecutive darts $g,g'$ of $\mathcal E(w)$. Whether
$g,g'$ lie in one $I_i$ or $g$ ends $I_i$, the whole slot lies in $S_i$
directly after $g$; as superblocks keep their linear order, every slot stays
a contiguous sequence with its old order, followed by a dart of $B$. Hence
$\rho^\tau_w$ restricts to every block rotation at $w$ and no two blocks
interleave. The resulting planar rotation system is $\rho^W$ except at $u,w$;
Lemma~\ref{lem:eng4-wheel}(b) gives $\rho^\tau$.

(b), (c) Rotations and anchors at vertices other than $u,w$ are unchanged,
and observations at $u$ occur only in (c), where the hypothesis applies; so
we consider observations at $w$. (b) A marked cycle $C'$ through $w$ avoids
$u$, so by
Lemma~\ref{lem:eng4-pole}(b) both its anchors at $w$ lie in one $I_i$ or in
one block $B'\ne B$, hence in one superblock $S_i$. For $g\in S_i$ the gap is
determined by the linear order of $S_i$; for $g\notin S_i$,
$g\in(\anout,\anin)$ if and only if $\anin$ precedes $\anout$ in $S_i$. Both
are preserved.

(c) By Lemma~\ref{lem:eng4-pole}(c), $\anin(C,w)\in S_a$ and $\anout(C,w)\in
S_b$ with $a\ne b$. A dart $g\in S_a$ lies in $(\anout(C,w),\anin(C,w))$ if
and only if it precedes $\anin(C,w)$ in $S_a$, and $g\in S_b$ if and only if
it follows $\anout(C,w)$ in $S_b$; these are preserved. For $g\in S_i$, $i\ne
a,b$: $g\in(\anout(C,w),\anin(C,w))$ if and only if $(S_b,S_i,S_a)$ is
counter-clockwise in the order of superblocks, if and only if $(e_a,e_i,e_b)$
is counter-clockwise at $u$ (the superblock order is the reverse of the
order of the $e_j$, for $\rho$ by (a) and for $\rho^\tau$ by definition). So
$\gap(C,w,g)=\gap(C,u,e_i)$ under both systems, equal by hypothesis.
\end{proof}

\begin{lemma}[bond map and orbits]\label{lem:eng4-orbits}
(a) If $\deg w=D$, let $\beta(e_i)$ be the unique dart of $I_i$
(Lemma~\ref{lem:eng4-pole}(d)). Then $\rho_w=\rev\beta(\rho_u)$ for every
$\rho$ as above, and $\rho^\tau_w=\rev\beta(\tau)$. If $u\in C$, then $\beta$
exchanges the anchor types of $C$ at $u$ and $w$, and
$\gap(C,u,e)=\gap(C,w,\beta e)$ whenever $\rho_w=\rev\beta(\rho_u)$.
(b) A permutation $\pi$ of a set of size $n$ fixes some cyclic order if and
only if all its orbits have equal length. If this holds, $n\ge3$ and $\pi$
preserves a pair $\{x,y\}$, then $\pi$ is the identity or a fixed-point-free
involution with $\pi(x)=y$. In the latter case every
$\pi$-invariant cyclic order $\sigma$ has exactly one of $z,\pi z$ in
$(x,y)_\sigma$ for each $z\ne x,y$; and if $\Lambda$ contains one element of
each orbit other than $\{x,y\}$ and $L$ is a linear order of $\Lambda$, then
$\sigma=(x,L,y,\pi L)$ is $\pi$-invariant with $(x,y)_\sigma=\Lambda$.
\end{lemma}

\begin{proof}
(a) Here $S_i=I_i=\{\beta e_i\}$; apply Lemma~\ref{lem:eng4-superblock}(a);
$C$ enters $w$ through $B_a$ (Lemma~\ref{lem:eng4-pole}(c)), and
Lemma~\ref{lem:v4-gaps}(b) applies to $\beta$. (b) $\pi$ fixes $\sigma$ if
and only if it commutes with the successor map of $\sigma$, an $n$-cycle,
that is, $\pi$ is a power of it; and a permutation with $n/k$ orbits of
length $k$ is conjugate to the $(n/k)$-th power of an $n$-cycle. Since
$\pi(x)\in\{x,y\}$, all orbits have length one, or all have length two with $\pi(x)=y$. A $\pi$-invariant $\sigma$ is shifted by $\pi$, for an
involution by half its length, so $y$ is opposite $x$ and $\pi$ maps
$(x,y)_\sigma$ onto $(y,x)_\sigma$. Finally $\pi(x,L,y,\pi L)=(y,\pi L,x,L)$.
\end{proof}

\begin{definition}[SimplifyMatching]\label{def:eng4-sm}
Let $u$ be as above with pipe $(u,r,\varphi)$, and assume that its partner
$w$ is unmatched (\emph{case (i)}), or $r=w$ (\emph{case (ii)}), or $w$ is a
block-vertex with trivial embedding tree matched by $(w,w',\varphi')$ with
$w'\ne u$ (\emph{case (iii)}). The graph part is that
of~\cite[Section~3.4]{BFR21}: delete the pipe of $u$ in cases (i), (ii),
and replace both pipes by $(r,w',\alpha)$, $\alpha=\varphi'\beta\varphi^{-1}$,
in case (iii). In case (ii) let $\pi=\varphi\beta^{-1}$.
\begin{enumerate}[label=(\alph*),nosep]
\item \emph{$u$ on no marked cycle:} in case (ii) reject unless all orbits of
  $\pi$ have equal length; no update.
\item \emph{$u\in C$, $r\ne w$:} add, for every nonanchor $e$ at $u$, the
  row $\langle C,u,e\rangle\oplus\langle C(r),r,\varphi e\rangle=
  \varepsilon(\varphi)$.
\item \emph{$u\in C$, $r=w$:} reject unless $\pi$ is the identity or a
  fixed-point-free involution exchanging $\anout(C,w)$ and $\anin(C,w)$; in
  the second case add, for every orbit $\{f,\pi f\}$ other than the anchor
  pair, the row $\langle C,w,f\rangle\oplus\langle C,w,\pi f\rangle=1$.
\end{enumerate}
\end{definition}

\begin{lemma}[SimplifyMatching]\label{lem:eng4-sm}
In the situation of Definition~\ref{def:eng4-sm}, $\deg w=D$ in cases (ii)
and (iii); the step satisfies the contract, yields an instance, and adds
at most $2D$ variables and $D$ rows.
\end{lemma}

\begin{proof}
In case (ii) the pipe gives $\deg w=D$, in case (iii)
Lemma~\ref{lem:eng4-pole}(d). In cases (i), (iii), $r\notin\{u,w\}$. Forward
maps keep $\rho$; every backward rotation system is $\rho'^\tau$ for some
$\tau$, so by Lemma~\ref{lem:eng4-superblock} it suffices to check the pipes and the
observations at $u,w$.
\emph{Off cycles (a).} If $\deg w=D$, a marked cycle through $w$ would avoid
$u$ and have both anchors in one single-dart set $I_i$
(Lemma~\ref{lem:eng4-pole}(b),(d)); so in cases (ii), (iii) $w$, and
by~\ref{it:v4-M4} $r$ and $w'$, lie on no marked cycle. Rejection in (ii) is
sound by Lemma~\ref{lem:eng4-orbits}(b), as a witness has $\pi(\rho_w)=\varphi\beta^{-1}(\rev\beta\rho_u)=
\rev\varphi(\rho_u)=\rho_w$. Forward, in case (iii)
$\rev\rho_{w'}=\varphi'(\rho_w)=\varphi'\rev\beta(\rho_u)=\alpha(\rho_r)$.
Backward: in cases (i), (iii) let $\tau=\rev\varphi^{-1}(\rho'_r)$; in case
(ii) let $\sigma$ be $\pi$-invariant and $\tau=\beta^{-1}(\rev\sigma)$. With
$\rho=\rho'^\tau$, the pipe of $u$ holds in (i), (iii) as $\rho_r=\rho'_r$;
in (ii) $\rho_w=\sigma$ and $\varphi(\tau)=\pi(\rev\sigma)=\rev\sigma$; in
(iii) $\varphi'(\rho_w)=\varphi'\rev\beta(\tau)=\alpha(\rho'_r)=
\rev\rho'_{w'}$. Observations at $w$ are kept
(Lemma~\ref{lem:eng4-superblock}(b)).
\emph{Remote aligned (b).} Here $w\in C$ (Lemma~\ref{lem:eng4-pole}(c)) and
$\varphi(e)$ is nonanchor at $r\in C(r)$. Forward,
Lemma~\ref{lem:v4-gaps}(b) gives the rows. Backward, with
$\tau=\rev\varphi^{-1}(\rho'_r)$, Lemma~\ref{lem:v4-gaps}(b) and the rows
give $\gap_\tau(C,u,e)=\gap_{\rho'}(C(r),r,\varphi e)\oplus
\varepsilon(\varphi)=\gap_{\rho'}(C,u,e)$, so $\rho'^\tau$ keeps every
observation (Lemma~\ref{lem:eng4-superblock}(c)); the pipes hold as in (a).
In case (iii), $\varphi^{-1}$, $\beta$, $\varphi'$ map anchor pairs onto
anchor pairs, so the shortcut is aligned.
\emph{Self-pipe (c).} $w\in C$, and $\varphi$, $\beta$ map the anchor pair
at $u$ onto that at $w$, so $\pi$ preserves it, and by Lemma~\ref{lem:eng4-orbits}(b) the
rejection is sound and the two remaining cases are exhaustive. By Lemma~\ref{lem:eng4-orbits}(a),
$\rev\varphi(\rho_u)=\pi(\rho_w)$, so the pipe holds if and only if
$\pi(\rho_w)=\rho_w$; for $\pi=\mathrm{id}$ it always holds. For an
involution, the rows hold forward by Lemma~\ref{lem:eng4-orbits}(b).
Backward, the set $\Lambda$ of nonanchor $f$ with $\gap_{\rho'}(C,w,f)=0$
contains one dart of each nonanchor orbit by the rows; let
$\sigma=(\anout(C,w),L,\anin(C,w),\pi L)$ for a linear order $L$ of
$\Lambda$ and $\tau=\beta^{-1}(\rev\sigma)$. Then $\sigma$ is
$\pi$-invariant, $\gap_\sigma(C,w,f)=\gap_{\rho'}(C,w,f)$, and by
Lemma~\ref{lem:eng4-orbits}(a) $\gap_\tau(C,u,e)=\gap_\sigma(C,w,\beta e)=
\gap_{\rho'}(C,u,e)$. So $\rho'^\tau$ keeps every observation and
$\rho'^\tau_w=\sigma$.
\emph{Invariant and size.} The graph and the marked cycles are unchanged;
the shortcut is off in (a) and aligned in (b); at most $2(D-2)$ observers
and $D-2$ rows are added.
\end{proof}

\begin{remark}[the self-pipe rows are necessary]\label{rem:eng4-selfpipe}
Let $G$ have two vertices $u,w$ and parallel edges $a,b,d,e$, the marked
digon $C:u\xrightarrow{a}w\xrightarrow{b}u$, observers $x_d,x_e$ of
$(C,w,d),(C,w,e)$, the row $x_d\oplus x_e=0$, and the aligned pipe
$(u,w,\varphi)$ exchanging the ends of $a,b$ and of $d,e$. Then
$\pi=(a\,b)(d\,e)$, so ordinary SimplifyMatching deletes the pipe, after
which $(a,d,e,b)$ at $w$ satisfies the row. With the pipe, $\rho_w$ is
$(a,d,b,e)$ or $(a,e,b,d)$, so $d,e$ lie in different gaps and the instance
is infeasible; the rows of Definition~\ref{def:eng4-sm}(c) record exactly
this. In the identity case nothing is lost, since $\rho_w=\rev\beta(\rho_u)$
already follows from planarity (Lemma~\ref{lem:eng4-orbits}(a)). We do not
know whether the involution case occurs in instances derived from compiled
instances.
\end{remark}

\subsection{The terminal solver}\label{sec:eng4-terminal}

Let $\mathcal I$ be a pipe-free instance with $G^W$ planar, and for a marked
cycle $C$ let $B(C)$ be the block of $G^W$ containing $C^W$. For a P-vertex
$v$ on a marked cycle we write $C=C(v)$. The \emph{terminal system}
$\Sigma(\mathcal I)$ consists of $R$ and the rows (R1)--(R3) below, over $X$
and the following new variables, each created only when a row uses it:
$\zeta_\mu$ for an R node $\mu$ of a block of $G^W$; $\xi_{\mu,\varepsilon}$
for a P node $\mu$ of $\mathcal T_{B(C)}$ that displays a lifted marked cycle
$C^W$ and a nonanchor direction $\varepsilon$ of $\mu$; and $\eta_{v,B'}$ for
a P-vertex $v$ on a marked cycle $C$ that is a cutvertex of $G^W$ and a block
$B'\ne B(C)$ of $G^W$ at $v$.
\begin{enumerate}[label=(R\arabic*),nosep]
\item\label{it:eng4-R1} For an observer $x$ of $(C,v,d)$ such that the edge
  of $d$ lies in a block $B_d\ne B(C)$: $x\oplus\eta_{v,B_d}=0$.
\item\label{it:eng4-R2} For an observer $x$ of $(C,v,d)$ with $d$ in $B(C)$,
  let $m$ be the median of the leaves $\anout(C,v),d,\anin(C,v)$ in
  $\Theta_{B(C)}(v)$, $\mu$ its source, and $e_{\mathrm{out}},e_d,
  e_{\mathrm{in}}$ the directions at $m$ towards them. If $\mu$ is an R node:
  $x\oplus\zeta_\mu=\operatorname{or}_{\theta}(e_{\mathrm{out}},e_d,
  e_{\mathrm{in}})$, $\theta$ the reference order of $m$. If $\mu$ is a P
  node: $x\oplus\xi_{\mu,e_d}=0$.
\item\label{it:eng4-R3} For every Q-vertex $q$ with $\deg q\ge3$:
  $f(q)\oplus\zeta_{\mu(q)}=b_q$, where $\mu(q)$ is the R node whose skeleton
  contains the wheel of $q$, and $b_q=0$ if and only if the reference
  embedding of $\mu(q)$ induces $\psi(q)$ at $q$ (spokes correspond to
  skeleton edges).
\end{enumerate}
In~\ref{it:eng4-R2} the source matters: a P node with three directions,
displayed as a Q node by the embedding tree of~\cite{BFR21}, gets a free
variable $\xi_{\mu,\varepsilon}$, never a reflection variable. A digon block
has no nonanchor dart, so observations at its vertices use~\ref{it:eng4-R1}.

\begin{theorem}[terminal solver]\label{thm:eng4-terminal}
Let $\mathcal I$ be a pipe-free instance of size $S$. If $G^W$ is not
planar, $F(\mathcal I)=\emptyset$. Otherwise $\Sigma(\mathcal I)$ consists of
$R$ and $O(S)$ further rows, all binary, over $X$ and $O(S)$ new variables;
it is computed in $O(S)$ time; $F(\mathcal I)$ is the projection to $X$ of
its solution set; and every solution yields in $O(S)$ time a witness of
$\mathcal I$ whose assignment is its restriction to $X$.
\end{theorem}

\begin{proof}
If $\mathcal I$ has a witness, $G^W$ is planar by
Lemma~\ref{lem:eng4-wheel}(a). Let $G^W$ be planar.
\emph{Well-definedness.} In~\ref{it:eng4-R2}, a P median displays $C^W$
(Lemma~\ref{lem:v4-frontier}(c)) and no other lifted marked cycle
(Lemma~\ref{lem:eng4-qint}(d)), so $\xi_{\mu,\varepsilon}$ refers to the
pole-independent gap $\gap_\mu$ with respect to $C^W$. In~\ref{it:eng4-R1},
$v$ is a cutvertex on the single marked cycle $C$. In~\ref{it:eng4-R3}, the
wheel of $q$ is $3$-connected and lies in the skeleton of one R node
$\mu(q)$~\cite[proof of Lemma~8]{BFR21}. A spoke $qp$ lies in an expansion
containing the skeleton vertices $q,p$, hence in that of the skeleton edge
$qp$; so every skeleton edge at $q$ contains exactly one spoke,
$\Theta_B(q)$ is the single node of $\mu(q)$, and $\mathcal E(q)$ is the
order of the spokes in $\mathcal E_{\mu(q)}(q)$ for every embedding
$\mathcal E$. Likewise a rim edge $pp'$ lies in the expansion of the
skeleton edge $pp'$. So the skeleton contains the wheel, with some rim
edges possibly represented by virtual edges, and as the wheel is
$3$-connected, the reference embedding induces $\psi(q)$ or $\rev\psi(q)$ at $q$.
\emph{Forward.} Given a witness $(\rho,\alpha)$, extend $\rho$ to $\rho^W$
(Lemma~\ref{lem:eng4-wheel}(a)) with block embeddings $\mathcal E$, and put
$\zeta_\mu=r_\mu(\mathcal E)$ and $\xi_{\mu,\varepsilon}=\gap_\mu(\mathcal
E,\varepsilon)$. The darts of $B'\ne B(C)$ at $v$ do not interleave with the
anchors of $C$, which lie in $B(C)$ (Lemma~\ref{lem:v4-assembly}(a)), so they
have one gap; let $\eta_{v,B'}$ be it. Then~\ref{it:eng4-R1} holds
by~\ref{it:v4-W3}, \ref{it:eng4-R2} by Lemmas~\ref{lem:eng4-wheel}(d)
and~\ref{lem:v4-frontier}(c), and~\ref{it:eng4-R3} by~\ref{it:v4-W2}.
\emph{Backward.} Given a solution, embed each R skeleton by its reference
embedding or the reverse according to $\zeta_\mu$, each P skeleton
displaying a lifted marked cycle (at most one) by an order
$(\varepsilon^{\mathrm{out}},L_0,\varepsilon^{\mathrm{in}},L_1)$ with the
directions of gap $0$ in $L_0$ and of gap $1$ in $L_1$, and everything not
determined by a created variable arbitrarily. At every cutvertex $x$ of
$G^W$ start from the rotation of one block and insert those of the others as
contiguous sequences; if $x$ is a P-vertex on a marked cycle $C$, start with
$B(C)$ and insert $B'$ directly after $\anout(C,x)$ if $\eta_{x,B'}=0$ and
after $\anin(C,x)$ otherwise. Block assembly
(Lemma~\ref{lem:v4-assembly}(c)) gives a planar $\rho^W$, and
Lemma~\ref{lem:eng4-wheel}(b) gives $\rho$. Later insertions keep earlier
blocks in their interval, so~\ref{it:eng4-R1} gives~\ref{it:v4-W3} for
$d\notin B(C)$; for $d\in B(C)$ it follows from~\ref{it:eng4-R2} and
Lemma~\ref{lem:v4-frontier}(c); \ref{it:eng4-R3} gives~\ref{it:v4-W2}.
\emph{Size and time.} Each observer and each Q-vertex gives one row and at
most one new variable. Building $G^W$, testing planarity, and computing
blocks, SPQR trees and reference embeddings takes
$O(S)$~\cite{HT74,Tar72,HT73,GM01}, and so does the complete family of
embedding trees with pointers from darts to leaves, skeleton edges of
directions and reference orders (Lemma~\ref{lem:v4-trees}). With cyclic
position indices in reference orders and in every $\psi(q)$, a triple is
oriented in $O(1)$. For~\ref{it:eng4-R2}, all observations of a P-vertex $v$
inside $B(C)$ use the same anchor leaves of $\Theta_{B(C)}(v)$; one traversal
of the path between them and of the subtrees hanging off it records, for
every leaf, its median and direction, in $O(|\Theta_{B(C)}(v)|)$, and each
tree is used for at most one cycle. \ref{it:eng4-R1} uses the block of each
dart, \ref{it:eng4-R3} the single node of $\Theta_B(q)$ and one triple. The
binary system is solved in linear time by breadth-first search (rows are
weighted edges, possibly loops; consistency means that every cycle has even
weight). For the witness, give every node of the complete family the
rotation induced by its source (Lemma~\ref{lem:v4-trees}), read every block
rotation as a frontier, and assemble with doubly linked cyclic lists, $O(1)$
splices per insertion; all in $O(S)$.
\end{proof}

\subsection{The engine}\label{sec:eng4-main}

The \emph{schedule} is that of the proof of~\cite[Lemma~7]{BFR21}: (S1) if a
pipe joins two cutvertices, apply EncapsulateAndJoin to it; (S2) otherwise,
if a matched block-vertex has a nontrivial embedding tree, apply PropagatePQ
at it; (S3) otherwise apply SimplifyMatching at $u$, where $u$ is the other
end of the pipe of a matched cutvertex of maximum degree if one exists, and
any matched vertex otherwise. By the proof of~\cite[Lemma~7]{BFR21}, if an
instance with all P-vertices of degree at least four and $G^W$ planar has a
pipe, the selected operation is applicable: in (S3), $u$ is a block-vertex
with trivial embedding tree, and its partner is unmatched or a matched
block-vertex with trivial embedding tree, which is
Definition~\ref{def:eng4-sm}.

\begin{algorithm}[t]
\caption{The local-gap engine}\label{alg:eng4}
\begin{algorithmic}[1]
\State make the input $\mathcal I$ canonical; apply ConvertSmall to all P-vertices of
  degree $\le3$; call the result $\mathcal I_0$; $j\gets0$
\While{$\mathcal I_j$ has a pipe}
  \If{$G_j^W$ is not planar} \Return reject \EndIf
  \State apply the step selected by the schedule
  (Lemmas~\ref{lem:eng4-join}, \ref{lem:eng4-ppq}, \ref{lem:eng4-sm});
  \textbf{if} it rejects, \Return reject
  \State apply ConvertSmall to all P-vertices of degree $\le3$; call the
  result $\mathcal I_{j+1}$; $j\gets j+1$
\EndWhile
\State solve $\Sigma(\mathcal I_j)$ (Theorem~\ref{thm:eng4-terminal});
  \textbf{if} it is inconsistent, \Return reject
\State construct a witness of $\mathcal I_j$; apply the backward maps in
  reverse order
\end{algorithmic}
\end{algorithm}

\begin{theorem}[engine]\label{thm:eng4-main}
Let $\mathcal I$ be an instance of size $S_0$, let $m_0=|E(G)|$, let $D_0$ be
the maximum degree of a pipe end ($D_0=0$ if there is no pipe), and
$S_\ast=S_0+m_0D_0$. Algorithm~\ref{alg:eng4} performs at most $2m_0$
iterations of the while loop, every intermediate instance has size
$O(S_\ast)$, and:
\begin{enumerate}[label=(\alph*),nosep]
\item if it rejects, $F(\mathcal I)=\emptyset$;
\item otherwise it outputs a system $\Sigma$ of binary rows over a set of
  variables containing $X$, consisting of $R$ and $O(S_\ast)$ further rows,
  such that $F(\mathcal I)$ is the projection to $X$ of the solution set of
  $\Sigma$, and from every solution it constructs a witness of $\mathcal I$
  whose assignment is the solution restricted to $X$.
\end{enumerate}
The running time, including the witness, is
$O\bigl((m_0+1)(S_0+m_0D_0)\bigr)$ word operations.
\end{theorem}

\begin{proof}
\emph{Correctness.} Every step satisfies the contract and produces an
instance (Section~\ref{sec:v4-lang}; Lemmas~\ref{lem:eng4-convert},
\ref{lem:eng4-join}, \ref{lem:eng4-ppq}, \ref{lem:eng4-sm}), and a rejection
at the planarity test is sound by Lemma~\ref{lem:eng4-wheel}(a). Composing
the contracts and the backward maps, with one assignment throughout, and
applying Theorem~\ref{thm:eng4-terminal} to the final instance gives (a) and
(b).
\emph{Termination.} Our trees keep their leaves as degree-two vertices
(so that darts keep their identities and marked cycles can be rerouted),
which become Q-vertices; the size analysis in the proof
of~\cite[Lemma~10]{BFR21} counts only inner tree nodes, identifying each
leaf with its reattached edge. Suppressing our leaves yields the graph
counted there, and subdividing an edge
changes neither the degree, the cut status nor the embedding tree of any
other vertex, nor partner poles. The schedule and the potential $\Phi$
of~\cite[Section~3.6]{BFR21} depend only on these data and on the pipes, not
on Q-cells, observations or rows, and ConvertSmall is applied as assumed
in~\cite{BFR21}; so our run makes the same choices with the same
potential. \cite[Lemma~10]{BFR21} accounts for the pipes an
operation removes and creates; every other pipe keeps its class, since its
ends keep their cut status (Lemma~\ref{lem:eng4-encap} for
EncapsulateAndJoin; contraction of trees avoiding $y$ turns $G'-y$ into
$G-y$ for PropagatePQ; the graph is unchanged otherwise). New pipes of degree
at most three are cap pipes at block-vertices (Lemma~\ref{lem:eng4-encap}),
of weight $0$. So each step decreases $\Phi$ by
\cite[Lemma~10(i)]{BFR21} (which uses~\cite[Lemma~9]{BFR21}), and
$\Phi(\mathcal I_0)\le2m_0$ by~\cite[Lemma~11]{BFR21}. Nothing else is
imported: the size bounds of~\cite[Lemma~10(ii),(iii) and Theorem~12]{BFR21}
do not cover our leaves, so we bound sizes separately.
\emph{Size.} New pipes have degree at most $D_0$; each step, including the
tree leaves and the following conversion, adds $O(D_0)$; canonicalization
and the initial conversion add $O(S_0)$; so all sizes are $O(S_\ast)$.
\emph{Time.} Let $S=O(S_\ast)$ bound all sizes. Canonicalization takes
$O(S_0)$, as an observable triple is determined by its dart. Per iteration,
$G^W$, planarity, blocks, SPQR trees and the selection take $O(S)$ (a matched
block-vertex has a trivial tree if and only if no R node contains it and
exactly one P node has it as a pole), and so does each step (for
EncapsulateAndJoin including the planarity test of $K$ and a breadth-first
search from $\{p,q\}$ for the targets; Lemma~\ref{lem:v4-trees} for
PropagatePQ; partner, sets $I_i$ and orbits for SimplifyMatching), the
conversion and each backward map
(Lemma~\ref{lem:eng4-stars}(c) for the join). With at most $2m_0$
iterations and the terminal stage in $O(S_\ast)$, the total is
$O(S_0+m_0S_\ast+S_\ast)=O\bigl((m_0+1)(S_0+m_0D_0)\bigr)$.
\end{proof}

\begin{remark}[why we restrict anchor pairs]\label{rem:eng4-anchors}
Arbitrary collections of local triple orientations at one P-vertex need not
admit an affine description. At a vertex with darts $a,b,c,d$ and
unconstrained rotation $\sigma$, the vector
$(\operatorname{or}_\sigma(a,b,c),\operatorname{or}_\sigma(a,b,d),
\operatorname{or}_\sigma(a,c,d),\operatorname{or}_\sigma(b,c,d))$ of local
gaps for several anchor pairs takes the six values $0000$, $0011$, $1001$,
$1100$, $0110$, $1111$ (for $\sigma=(a,b,c,d)$, $(a,b,d,c)$, $(a,c,b,d)$,
$(a,c,d,b)$, $(a,d,b,c)$, $(a,d,c,b)$), while the projection of the solution
set of a system of rows, even with auxiliary variables, is an affine space of
size $0$ or a power of two; compare~\cite[Section~8.1]{Pfr22}. This excludes
affine descriptions only, not other graph gadgets; and different anchor pairs
are not by themselves the obstruction, since $\operatorname{or}_\sigma(1,3,2)$
and $\operatorname{or}_\sigma(1,4,3)$ take all four values on darts
$1,2,3,4$. Theorem~\ref{thm:eng4-main} avoids the question: all observations
at a P-vertex use one anchor pair, for which every vector of gaps is
realizable.
\end{remark}
